\section{Conclusion}

Under a fully matched training budget, a two-block matrix fast-weight
model stores and retrieves in-context bindings that neither a
parameter-matched flat-vector recurrence nor a parameter-matched
transformer learns at all, and it does so from a state whose size does not
grow with context: recall stays at ceiling to eight times the binding
span % <!-- evidence: C3 -->
at a constant 32{,}768
bytes, % <!-- evidence: C10 -->
while no KV-cache budget up to thirty-two
times those bytes changes the baseline's chance-level
reading. % <!-- evidence: C4 -->
The
bindings are causally localized to the first block's fast-weight matrix
and stored in a form only the model's own forward pass decodes. Because
the baseline never becomes competent, the comparison yields no memory
multiplier, and we report it under its registered phrasing instead:
baseline non-competitive at matched params/tokens. For the workshop's
question, whether long-horizon reasoning must pay a memory cost that grows
with context, this is a small-scale existence result on the recall
substrate such reasoning requires: at matched budget, the constant-memory
architecture was the only one that learned to remember.
